\section{How many transversals? Semiqueens and queens}\label{sec:counting}

\topic{Counting complete mappings}

The first result on the number of transversals concerned cyclic groups and was conjectured by Vardi.

\begin{theorem}[Eberhard--Manners--Mrazovi\'c~\cite{emm19}; Vardi's conjecture]\label{thm:emm19}
For odd $n$, the Cayley table of $\Z_n$ has
\[
\bigl(e^{-1/2}+o(1)\bigr)\,n\,\frac{n!^2}{n^n}
\]
transversals.
\end{theorem}

Equivalently, this counts the ways to place $n$ non-attacking \term{semiqueens} on an $n\times n$ board wrapped into a torus; a semiqueen moves along rows, columns and one family of diagonals. Vardi~\cite{var91} conjectured in 1991 that the count lies between $c_1^n\,n!$ and $c_2^n\,n!$ for constants $0<c_1\leq c_2<1$, the lower bound being the hard part. Theorem~\ref{thm:emm} extends Theorem~\ref{thm:emm19} to every group, with the same constant $e^{-1/2}$ and an extra factor $|G^{ab}|$, and has the following consequence.

\begin{corollary}[Eberhard--Manners--Mrazovi\'c~\cite{emm22}]\label{cor:most}
For large $k$, among all groups of order $2^k$, the elementary abelian group $\Z_2^k$ has the most complete mappings.
\end{corollary}

\topic{The $n$-queens problem}

The semiqueens are a relative of a much older puzzle.

\begin{problem}[Bezzel~\cite{bez1848}, 1848]\label{prob:queens}
Place eight queens on a chessboard so that no two attack each other, that is, no two share a row, a column or a diagonal.
\end{problem}

Nauck~\cite{nau1850} found all $92$ solutions in 1850; Gauss also worked on the puzzle. Pauls~\cite{pau1874} showed in 1874 that on an $n\times n$ board $n$ non-attacking queens fit for every $n\notin\{2,3\}$. Let $Q(n)$ denote the number of solutions, so that $Q(8)=92$; the question is how $Q(n)$ grows.

\begin{figure}[ht]
\centering
\begin{tikzpicture}[scale=0.42]
    \foreach \i in {0,...,7}{\foreach \j in {0,...,7}{%
      \pgfmathtruncatemacro{\pty}{mod(\i+\j,2)}%
      \ifnum\pty=1 \fill[black!12] (\j,-\i) rectangle ++(1,-1);\fi}}
    \draw (0,0) rectangle (8,-8);
    \foreach \c [count=\r from 0] in {1,5,8,6,3,7,2,4}
      \fill[BlkConjecture] (\c-0.5,-\r-0.5) circle (0.3);
  \end{tikzpicture}

\caption{A solution of the eight-queens puzzle.}
\label{fig:eight-queens}
\end{figure}

Number the rows and the columns $1,\dots,n$. With one queen in each row and each column, the queen of row $j$ stands in column $\sigma(j)$ for a permutation $\sigma$, and the square $(j,c)$ lies on the diagonals $c+j=\text{const}$ and $c-j=\text{const}$.

\begin{fact}\label{fact:queens}
A permutation $\sigma$ is a solution if and only if the $n$ numbers $\sigma(j)+j$ are distinct and the $n$ numbers $\sigma(j)-j$ are distinct.
\end{fact}

For the solution in Figure~\ref{fig:eight-queens}:
\begin{center}
\renewcommand{\arraystretch}{1.2}
\begin{tabular}{r|cccccccc}
$j$ & 1&2&3&4&5&6&7&8\\\hline
$\sigma(j)$ & 1&5&8&6&3&7&2&4\\
$\sigma(j)+j$ & 2&7&11&10&8&13&9&12\\
$\sigma(j)-j$ & 0&3&5&2&$-2$&1&$-5$&$-4$
\end{tabular}
\end{center}

\topic{P\'olya's toroidal board}

P\'olya~\cite{pol1918} studied in 1918 the \emph{doubly periodic} solutions: those that remain non-attacking when the board is repeated to fill the whole plane (Figure~\ref{fig:queens-tiled}). Equivalently the diagonals wrap around, the board is a torus, and the numbers $\sigma(j)+j$ and $\sigma(j)-j$ must be distinct modulo $n$. The torus is the more regular problem, with $n$ rows, $n$ columns and $n$ diagonals in each direction. A toroidal solution is also an ordinary one, so the number $T(n)$ of toroidal solutions satisfies $T(n)\leq Q(n)$, and P\'olya showed that $T(n)>0$ if and only if $n\equiv1,5\pmod6$.

\begin{figure}[ht]
\centering
\begin{tikzpicture}[scale=0.3]
    \foreach \bx in {0,1}{\foreach \by in {0,1}{%
      \foreach \i in {0,...,4}{\foreach \j in {0,...,4}{%
        \pgfmathtruncatemacro{\pty}{mod(\i+\j,2)}%
        \ifnum\pty=0 \fill[black!8] (5*\bx+\j,-5*\by-\i) rectangle ++(1,-1);\fi}}
      \foreach \j in {0,...,4}{%
        \pgfmathtruncatemacro{\c}{mod(2*\j,5)}%
        \fill[BlkConjecture] (5*\bx+\c+0.5,-5*\by-\j-0.5) circle (0.32);}}}
    \draw[step=1, black!30] (0,0) grid (10,-10);
    \draw[very thick] (0,0) rectangle (10,-10);
    \draw[very thick] (5,0) -- (5,-10) (0,-5) -- (10,-5);
  \end{tikzpicture}

\caption{Five queens on the $5\times5$ torus, repeated periodically: still no attacks.}
\label{fig:queens-tiled}
\end{figure}

On the torus, both puzzles are statements about the group $\Z_n$:
\begin{center}
\renewcommand{\arraystretch}{1.4}
\begin{tabular}{lll}
\textbf{on the $n\times n$ torus} & \textbf{condition on $\sigma$ (mod $n$)} & \textbf{in the group $\Z_n$}\\\hline
$n$ semiqueens & $\sigma(j)+j$ distinct & a complete mapping\\
$n$ queens & $\sigma(j)+j$ and $\sigma(j)-j$ distinct & a strong complete mapping
\end{tabular}
\end{center}
Semiqueens exist if and only if $n$ is odd, which is the Hall--Paige conjecture for $\Z_n$; queens exist if and only if $n\equiv1,5\pmod6$, by P\'olya's theorem. In Figure~\ref{fig:toroidal-queens}, where $\sigma(j)=2j$ on $\Z_5$, both $\sigma(j)+j=3j$ and $\sigma(j)-j=j$ are bijections of $\Z_5$.

\begin{figure}[ht]
\centering
\begin{tikzpicture}[scale=0.62]
    \foreach \i in {0,...,4}{\foreach \j in {0,...,4}{%
      \pgfmathtruncatemacro{\pty}{mod(\i+\j,2)}%
      \ifnum\pty=0 \fill[black!8] (\j,-\i) rectangle ++(1,-1);\fi}}
    \draw (0,0) grid (5,-5);
    \foreach \j in {0,...,4}{%
      \pgfmathtruncatemacro{\c}{mod(2*\j,5)}%
      \fill[BlkConjecture] (\c+0.5,-\j-0.5) circle (0.3);
      \node[font=\scriptsize, text=CiteColor] at (-0.4,-\j-0.5) {$\j$};
      \node[font=\scriptsize, text=CiteColor] at (\j+0.5,0.4) {$\j$};}
  \end{tikzpicture}

\caption{Five non-attacking queens on the $5\times5$ torus: $\sigma(j)=2j$ on $\Z_5$.}
\label{fig:toroidal-queens}
\end{figure}

\topic{Counting queens}

Rivin, Vardi and Zimmermann~\cite{rvz94} conjectured in 1994 that $\log Q(n)$ and $\log T(n)$ are both $\Theta(n\log n)$. Luria~\cite{lur17} showed that $T(n)\leq\bigl((1+o(1))\,ne^{-3}\bigr)^n$ and conjectured that this is tight. Bowtell and Keevash~\cite{bk21} proved that $T(n)=\bigl((1+o(1))\,ne^{-3}\bigr)^n$ for $n\equiv1,5\pmod6$, and Luria and Simkin~\cite{ls21} and Bowtell and Keevash~\cite{bk21} showed that $Q(n)\geq\bigl((1-o(1))\,ne^{-3}\bigr)^n$. Finally Simkin~\cite{sim21} determined the ordinary count:
\[
Q(n)=\bigl((1\pm o(1))\,ne^{-\alpha}\bigr)^n,\qquad \alpha=1.942\pm0.003,
\]
which is about $(0.143\,n)^n$. The tools are those of the Hall--Paige problem: a random greedy process, followed by absorption.

Two further results concern completion and symmetry.

\begin{theorem}[Glock--Munh\'a Correia--Sudakov~\cite{gms22}]\label{thm:completion}
For large $n$, any $n/60$ non-attacking queens on the $n\times n$ board can be completed to a solution, while some $\approx n/4$ non-attacking queens cannot.
\end{theorem}

The completion question goes back to Nauck~\cite{nau1850}, and the proof goes through rainbow matchings in bipartite graphs.

\begin{theorem}[Dai--Kelly~\cite{dk24}]\label{thm:reflecting}
For all large $n$ there is a reflecting $n$-queens configuration, answering questions of Klarner (1967) and Slater.
\end{theorem}

The group-theoretic generalisation of toroidal queens, strong complete mappings, is the subject of Problem~\ref{prob:strong}.
